\subsubsection{FFT (iterativo, complex)}
\begin{lstlisting}[style=mystyle, language=C++]
// TC: O(N log N)
// usa: multiplicacion de polinomios con coeficientes double
#include <bits/stdc++.h>
using namespace std;

using cd = complex<double>;
const double PI = acos(-1);

void fft(vector<cd>& a, bool invert){
	int n = a.size();
	// bit-reversal permutation
	for(int i=1,j=0;i<n;i++){
		int bit = n >> 1;
		for(; j&bit; bit>>=1) j ^= bit;
		j ^= bit;
		if(i < j) swap(a[i], a[j]);
	}
	for(int len=2; len<=n; len<<=1){
		double ang = 2 * PI / len * (invert ? -1 : 1);
		cd wlen(cos(ang), sin(ang));
		for(int i=0; i<n; i+=len){
			cd w(1);
			for(int j=0; j<len/2; j++){
				cd u = a[i+j], v = a[i+j+len/2] * w;
				a[i+j] = u + v;
				a[i+j+len/2] = u - v;
				w *= wlen;
			}
		}
	}
	if(invert){
		for(cd& x : a) x /= n;
	}
}

vector<long long> multiplyFFT(const vector<long long>& a, const vector<long long>& b){
	vector<cd> fa(a.begin(), a.end()), fb(b.begin(), b.end());
	int n = 1;
	while(n < (int)a.size() + (int)b.size()) n <<= 1;
	fa.resize(n); fb.resize(n);
	fft(fa, false); fft(fb, false);
	for(int i=0;i<n;i++) fa[i] *= fb[i];
	fft(fa, true);
	vector<long long> res(n);
	for(int i=0;i<n;i++) res[i] = (long long)llround(fa[i].real());
	return res;
}

int main(){
	// multiply (1 + x + 2x^2) * (3 + 4x) = 3 + 7x + 10x^2 + 8x^3
	auto res = multiplyFFT({1,1,2}, {3,4});
	for(long long x : res) cout << x << " ";
	cout << "\n";
	return 0;
}
\end{lstlisting}
